% SPDX-License-Identifier: Apache-2.0
\section{Curriculare Level: Erwerbsordnung und Bewältigbarkeit}
\label{sec:levels}
\begin{principle}{Status: optionales Planungsprofil}
Die folgenden Bedingungen definieren einen curricularen Levelplan.
Sie ändern weder Grundgraphgültigkeit noch aktuelle Frontier oder
gespeicherte Mastery. Es wird damit kein bereits implementierter
Levelvalidator behauptet. Personalisierungsstufen, Schwierigkeit und
QA-Reifegrade sind andere Begriffe.
\end{principle}

Ein curricularer Level ist eine kontextgebundene Erwerbsetappe, etwa ein
Schuljahr, ein Semester oder eine fachübergreifende Stufe. Das anschauliche
Bild eines Wissensgebäudes hilft: Kompetenzen sind Bauteile, Voraussetzungen
tragende Beziehungen, ein Level eine geplante Bauetappe. Das Bild darf aber
nicht suggerieren, alle Fächer müssten vollständig abgeschlossen sein,
bevor irgendwo ein Ziel des nächsten Levels verfügbar wird.

\subsection{Plan, Eintrittsbasis und eindeutiger Ersterwerb}
Für einen aufgelösten Plan $\pi$ seien
\begin{align}
 T_\pi&\subseteq A &&\text{geplante Erwerbsziele},\\
 B_\pi&\subseteq A\setminus T_\pi &&\text{ausdrücklich angenommene Eintrittsbasis},\\
 n_\pi&\in\N_{>0},\qquad
 \ell_\pi:T_\pi\to\{1,\ldots,n_\pi\} &&\text{Erwerbslevel}.
\end{align}
Die Levelmengen und ihre Präfixe sind
\[
 L_i=\{g\in T_\pi\mid\ell_\pi(g)=i\},\quad
 L_{<i}=\bigcup_{j<i}L_j,\quad L_{\leq i}=\bigcup_{j\leq i}L_j.
\]
Jedes Ziel erhält in diesem Plan genau einen Ersterwerbslevel. Ein anderer
Bildungsgang darf es anders platzieren. Spätere Wiederholung, Vertiefung oder
Anwendung sind zusätzliche Aktivitäten, keine neue Zielidentität.
$B_\pi$ ist eine Planungsannahme; sie schreibt keine Lernergebnisse.
Persönliche Lücken in dieser Basis brauchen eine eigene Diagnose und
zusätzliche Lernzeit.

\begin{figure}[htbp]
\centering
\begin{tikzpicture}[node distance=10mm]
\node[goal,text width=24mm] (b) {Eintrittsbasis\\$B_\pi$};
\node[goal,text width=24mm,right=12mm of b] (l1) {Level 1\\$L_1$};
\node[goal,text width=24mm,right=12mm of l1] (l2) {Level 2\\$L_2$};
\node[goal,text width=24mm,right=12mm of l2] (l3) {Level 3\\$L_3$};
\draw[-{Stealth},draw=muted] (b)--(l1);
\draw[-{Stealth},draw=muted] (l1)--(l2);
\draw[-{Stealth},draw=muted] (l2)--(l3);
\coordinate (levelscenter) at ($(b.south)!0.5!(l3.south)$);
\node[info,anchor=north,text width=130mm] at ($(levelscenter)+(0,-7mm)$)
 {Jedes Level besitzt einen global gültigen Erwerbsweg\\und ein begründetes Zeitbudget.};
\end{tikzpicture}
\caption[Curriculare Erwerbsstufen]{Planungsreihenfolge, keine \code{requires}-Kanten zwischen
Levelcontainern. Einzelziele werden über ihre eigenen Voraussetzungen verfügbar.}
\label{fig:levels}
\end{figure}

\subsection{Strukturelle Machbarkeit durch einen Abschlusszeugen}
Ein Level $L_i$ ist strukturell machbar, wenn sich seine Ziele als
$a_1,\ldots,a_m$ aufzählen lassen, sodass
\begin{equation}
 a_j\in\Frontier\!\left(B_\pi\cup L_{<i}\cup
                         \{a_1,\ldots,a_{j-1}\}\right)
 \qquad(1\leq j\leq m).
 \label{eq:level-witness}
\end{equation}
Hier gilt die \emph{globale} Frontier. Vererbte und entlang offener Ziele verfolgte,
clusterbezogene und nach Referenzauflösung fachübergreifende Anforderungen
bleiben erhalten. Ein Normalmodusfilter würde gerade jene Anforderungen
ausblenden können, deren Einhaltung die Levelprüfung nachweisen soll.

Zur Diagnose dienen die benötigten Atome
\begin{equation}
 \ReqAtoms(g)=\bigcup_{\substack{p\in G\\(g,p)\in\Reff^+}}\Atoms(p).
 \label{eq:reqatoms}
\end{equation}
Aus einem Abschlusszeugen folgt notwendig
\begin{equation}
 \forall g\in T_\pi:\quad
 \ReqAtoms(g)\subseteq B_\pi\cup L_{\leq\ell_\pi(g)}.
 \label{eq:level-closure}
\end{equation}
Eine benötigte Kompetenz außerhalb $B_\pi\cup T_\pi$ ist ungeplant;
ein benötigtes Ziel mit höherem Erwerbslevel ist ein Reihenfolgefehler.
Voraussetzungen innerhalb desselben Levels sind erlaubt. Die Mengenbedingung
allein genügt bei Clusteranforderungen noch nicht: Der Selbstblocker aus
Abschnitt~\ref{sec:validity} erfüllt sie für $L_1=\{a\}$,
besitzt aber keinen Abschlusszeugen.

Ein endlicher Prüfalgorithmus startet für jedes Level mit
$M^{(0)}=B_\pi\cup L_{<i}$ und wiederholt
\begin{equation}
 M^{(r+1)}=M^{(r)}\cup\bigl(L_i\cap\Frontier(M^{(r)})\bigr).
\end{equation}
Er ist erfolgreich genau dann, wenn sein Fixpunkt ganz $L_i$ enthält.
Die Verfügbarkeit eines noch offenen Ziels kann durch zusätzliche
Erfüllungen nicht verschlechtert werden. Deshalb kann jede gleichzeitig
hinzugefügte Menge auch nacheinander erworben werden. Ein vorhandener
Abschlusszeuge kann umgekehrt nicht vorzeitig an diesem Verfahren scheitern.
Das Verfahren ist ein \emph{hypothetischer} Lerndurchlauf; es ändert keine
persönlichen Daten. Im rein atomaren DAG liefert eine voraussetzungsgetreue
topologische Ordnung den bekannten Spezialfall.

\subsection{Aufwand, Reserve und Budget}
Strukturelle Machbarkeit beweist noch nicht, dass ein Schuljahr zeitlich
ausreicht. Jede Zeitbehauptung nennt Referenzlernende, Eintrittsbedingungen,
Ziel- und Wahlauswahl, Kalender, Schätzmethode und Unsicherheit.
Alle Größen müssen in derselben kalibrierten Zeiteinheit vorliegen:
\begin{equation}
 \Work_\pi(i;\Theta_\pi)+\Reserve_\pi(i)\leq \Budget_\pi(i),
 \qquad i=1,\ldots,n_\pi.
 \label{eq:budget}
\end{equation}
Dabei sind $\Work\geq0$, $\Reserve\geq0$ und $\Budget>0$.
Reserve ist geschützt und wird nicht gleichzeitig als verplanbare Zeit
ausgegeben. Ein Nullpuffer ist möglich, aber offenzulegen. Bereits im
Arbeitsaufwand enthaltene Prüfungsvorbereitung oder Wiederholung darf nicht
ein zweites Mal als Reservebedarf gezählt werden.

Ein erklärtes Aktivitätenmodell kann
$\Work_\pi(i;\Theta_\pi)=\sum_{h\in\mathcal H_i}\widehat t_h(\Theta_\pi)$
verwenden. $\mathcal H_i$ enthält jeden tatsächlichen Aufwand genau einmal:
Ersterwerb, gezielte Übung, Abruf, Prüfung und weitere geplante Arbeit.
Dasselbe kanonische Atom wird nicht in Mathematik und Physik zweimal als
Ersterwerb gezählt; eine zusätzliche Anwendung in Physik kann dennoch
eigenständige Zeit kosten. Unkalibrierte \glqq Zielpunkte\grqq{} oder
\code{weight}-Werte sind keine Minutenschätzung.
Tabelle~\ref{tab:workload} veranschaulicht die getrennte Verbuchung von Arbeit und Reserve.

\begin{table}[htbp]
\centering
\begin{tabular}{@{}lrr@{}}
\toprule
Illustrative Planung & Stunden & Bedeutung\\
\midrule
Ersterwerb & 42 & Bereits um gemeinsame Ziele bereinigt\\
Übung und Anwendung & 12 & Zusätzliche Tätigkeiten\\
Prüfung und Wiederholung & 6 & Im Arbeitsaufwand enthalten\\
Geschützte Reserve & 10 & Nicht für obige Arbeit verplant\\
Verfügbare Zeit & 75 & Nach Kalenderausschlüssen\\
\bottomrule
\end{tabular}
\caption{Rechenbeispiel: $\Work=60$ und $\Reserve=10$ passen in
$\Budget=75$. Die erfundenen Zahlen demonstrieren nur die Rechnung,
keine empirisch abgesicherte Dauer eines SkillPilot-Curriculums.}
\label{tab:workload}
\end{table}

Ein Gesamtbudget genügt nicht, wenn Fächer, Räume oder Termine eigene
bindende Grenzen haben. Für eine zusätzlich deklarierte Fächerzuordnung
müssen auch die jeweiligen $\Work_{i,s}+\Reserve_{i,s}\leq \Budget_{i,s}$
gelten, ohne Doppelzählung gemeinsam zugeordneter Zeit. Ein Kalenderplan
muss außerdem Voraussetzungen vor der konkreten Anwendung bereitstellen.

\subsection{Gültigkeitsanspruch und Abschlussgrenzen}
Ein Levelplan ist unter diesem optionalen Profil gültig, wenn er sowohl
strukturell als auch unter seinen erklärten Annahmen zeitlich machbar ist.
Ohne belastbare Aufwandsschätzung ist allein die strukturelle Aussage
zulässig. Der Anspruch bindet Plan, Curriculumfassung, aufgelösten Umfang
und Schätzannahmen. Er garantiert keine identische Lerngeschwindigkeit
für jede Person.

Der Abschluss eines Levels wird aus seiner vollständigen Zielmenge und den
Zieltypverträgen abgeleitet. Er ist keine pauschale Voraussetzung für alle
Ziele des nächsten Levels. Soll ein Teil der Arbeit durch Prüfungsnachweise
ersetzt werden, benötigt das den folgenden gesonderten Abschlussvertrag.

\paragraph{Quellen.}
\srcref{docs/concept/skill-graph/curriculum-levels.md}{Curriculum Levels};
\srcref{docs/concept/skill-graph/graph-definition.md}{Graph Definition, Abschnitt 12};
\srcref{docs/concept/didactic/curriculum-time-axis-and-pacing.md}{Zeitachse und Kapazität}.
